\section{Integrable coefficients and the controlled diffusion}
\label{sec:regularity}

The sign inequalities will first be proved for finite step functions.  We
therefore need estimates which survive approximation of an arbitrary
$\gamma\in\mathcal U$, including coefficients which diverge at time one.
All estimates in this section are on strips ending strictly before one.
The control and heat-semigroup methods follow the Parisi PDE framework of
Auffinger and Chen~\cite{AC2015}, Jagannath and Tobasco~\cite{JT2016},
and the zero-temperature treatment of Chen, Handschy, and
Lerman~\cite[Proposition~2 and Theorem~5]{CHL2018}.
We include the proofs to make the approximation uniformity explicit.
We retain the notation $u_\gamma=\partial_x\Phi_\gamma$ from the model
definition, and write $P_t$ for the heat semigroup with variance $t$.

\begin{lemma}[Compact-strip regularity and stability]
\label{lem:regularity}
For every $\gamma\in\mathcal U$, the function $\Phi_\gamma$ is even,
convex, and one-Lipschitz in its spatial variable.  It satisfies
\begin{equation}
 |x|\leq\Phi_\gamma(t,x)
 \leq |x|+\sqrt{\frac{2(1-t)}\pi}
               +\frac12\int_t^1\gamma(s)\,\dd s.
 \label{eq:terminal-uniform}
\end{equation}
For each integer $k\geq1$ and $T<1$, the derivative
$\partial_x^k\Phi_\gamma$ is bounded and jointly continuous on
$[0,T]\times\R$.  These bounds are uniform when the coefficients have
uniformly bounded $L^1(0,1)$ norms.  The derivatives are locally
absolutely continuous in time and satisfy the spatially differentiated
Parisi equation for almost every time.

If $\gamma_n\to\gamma$ in $L^1(0,1)$, then, for every $k\geq1$ and
$T<1$,
\begin{equation}
 \sup_{[0,T]\times\R}
 |\partial_x^k\Phi_{\gamma_n}-\partial_x^k\Phi_\gamma|
 \longrightarrow0.
 \label{eq:derivative-convergence}
\end{equation}
The same conclusion holds when each $\gamma_n$ is a finite step
function and its terminal datum is
$M_n^{-1}\log(2\cosh(M_nx))$, where $M_n\to\infty$.
\end{lemma}

\begin{proof}
For a nonnegative finite step coefficient and a smooth even convex
one-Lipschitz terminal datum $\psi$, backward evolution across an
interval with coefficient $c$ is
\begin{equation}
 \Phi(t,x)=
 \begin{cases}
 c^{-1}\log P_{b-t}(e^{c\Phi(b,\cdot)})(x),&c>0,\\
 P_{b-t}\Phi(b,\cdot)(x),&c=0.
 \end{cases}
 \label{eq:constant-heat}
\end{equation}
These formulas preserve evenness, convexity, and the one-Lipschitz
bound.  Convexity in the first case follows either from H\"older's
inequality applied after a spatial translation, or by differentiating
the logarithm: its second derivative is the tilted mean of the old
second derivative plus $c$ times the tilted variance of the old first
derivative.  The first derivative is a tilted mean of the old first
derivative.  The assertions in the case $c=0$ follow directly from
convolution.

For these step solutions, It\^o's formula and completion of the square
give the control representation
\begin{equation}
 \Phi(t,x)=\sup_{|a|\leq1}\E\left[
 \psi\left(x+B_1-B_t+\int_t^1\gamma(s)a_s\,\dd s\right)
 -\frac12\int_t^1\gamma(s)a_s^2\,\dd s\right].
 \label{eq:control-representation}
\end{equation}
The supremum is over progressively measurable controls.  Indeed, the
drift of $\Phi(s,Y_s)-\frac12\int_t^s\gamma a^2$ is
$-\gamma(a-u_\gamma)^2/2$ when
$\dd Y_s=\gamma(s)a_s\,\dd s+\dd B_s$.  Equality is attained by
$a_s=u_\gamma(s,Y_s)$; for the smooth finite-step problem the feedback
equation has a unique strong solution.  Consequently two such
solutions obey
\begin{equation}
 \|\Phi_{\gamma,\psi}-\Phi_{\widetilde\gamma,\widetilde\psi}
       \|_\infty
 \leq\|\psi-\widetilde\psi\|_\infty
             +\frac32\|\gamma-\widetilde\gamma\|_1.
 \label{eq:value-stability}
\end{equation}
Here the supremum is over the whole time-space domain, even though the
individual functions grow linearly.  The bound follows by using the
same controls on both sides of \eqref{eq:control-representation}.

Nonnegative nondecreasing finite step functions are dense in
$\mathcal U$ for the $L^1$ norm.  For example, truncate at a time
increasing to one, approximate the resulting bounded monotone
function by monotone step functions, and then let the truncation time
tend to one.  Since
$0\leq M^{-1}\log(2\cosh(Mx))-|x|\leq M^{-1}\log2$,
\eqref{eq:value-stability} defines a unique uniform limit, with control
representation \eqref{eq:control-representation} and terminal datum
$|x|$.  This agrees with the Parisi solution in the model definition.
All the asserted shape properties pass to the limit.  The zero control
and Jensen's inequality give the lower bound in
\eqref{eq:terminal-uniform}.  For the upper bound use the triangle
inequality and $|a|-a^2/2\leq1/2$ in the control representation.

We give the smoothing estimates explicitly.  This also establishes
the uniformity needed below.  On any strip $[0,b]$ with $b<1$,
monotonicity gives
\begin{equation}
 \gamma(t)\leq\frac{\|\gamma\|_1}{1-b},\qquad 0\leq t\leq b.
 \label{eq:local-coefficient-bound}
\end{equation}
Thus all coefficients in an $L^1$-bounded family have a common bound
$\Gamma$ there.  Work first with the smooth-terminal step solutions
and put $U=\Phi_x$.  Differentiation of the mild equation gives
\begin{equation}
 U(t)=P_{b-t}U(b)
       +\int_t^b\gamma(s)P_{s-t}[U(s)U_x(s)]\,\dd s.
 \label{eq:burgers-mild}
\end{equation}
We always have $\|U\|_\infty\leq1$, independently of its terminal
smoothing.  The heat-kernel estimate
$\|\partial_xP_h f\|_\infty\leq C h^{-1/2}\|f\|_\infty$
implies
\begin{equation}
 \|U_x(t)\|_\infty
 \leq \frac{C}{\sqrt{b-t}}
       +C\Gamma\int_t^b
                  \frac{\|U_x(s)\|_\infty}{\sqrt{s-t}}\,\dd s.
 \label{eq:fractional-bound}
\end{equation}
The constants do not involve derivatives of the terminal datum.

For completeness, the integral inequality in
\eqref{eq:fractional-bound} gives a uniform bound away from $b$.
The convolution of $h^{-1/2}$ with itself equals $\pi$.
Substituting the inequality once into its integral therefore reduces
it to an ordinary backward Gronwall inequality, with forcing bounded
by a constant times $1+(b-t)^{-1/2}$.  That forcing is integrable.
The resulting estimate is uniform on every smaller strip.

The same argument bounds all higher derivatives.  Suppose that
derivatives of $U$ through order $m-1$ are already bounded on a
slightly larger strip.  Differentiate \eqref{eq:burgers-mild} $m$
times, placing one derivative on the heat kernel and the other
$m-1$ on $UU_x$.  Its highest-order term is
$U\partial_x^mU$; every remaining product involves derivatives of
orders at most $m-1$.  The initial term is bounded by
$C(b-t)^{-1/2}\|\partial_x^{m-1}U(b)\|_\infty$.
Thus the same integral inequality applies to
$\|\partial_x^mU(t)\|_\infty$, with an additional bounded forcing.
Choose successively smaller strip endpoints in this induction.
For every fixed $m$ this proves the asserted bound, uniformly for
the approximants.

Uniform convergence of the functions now gives convergence of all
their spatial derivatives without a separate singular-kernel
comparison.  If a bounded $C^2$ function $F$ on $\R$ has bounded
second derivative, Taylor's formula gives, for every $h>0$,
\begin{equation}
 \|F'\|_\infty\leq
 \frac{2\|F\|_\infty}{h}
       +\frac h2\|F''\|_\infty.
 \label{eq:interpolation}
\end{equation}
Apply this to a difference of two approximating solutions, uniformly
on a fixed smaller strip.  The uniform second-derivative bound and
uniform convergence imply convergence of the first derivatives.
Apply \eqref{eq:interpolation} successively to the differences of
the first, second, and later derivatives; the next two derivative
bounds are available at every stage.  This proves uniform convergence
of every positive spatial derivative.  Their limits are the spatial
derivatives of the limiting function, by the fundamental theorem of
calculus.  Since the approximating derivatives are jointly continuous,
so are their uniform limits.

The argument applies also to an arbitrary $L^1$-convergent sequence in
$\mathcal U$: first obtain these bounds for each member by step
approximation, then use \eqref{eq:value-stability}, the uniform local
coefficient bound \eqref{eq:local-coefficient-bound}, and
\eqref{eq:interpolation}.  Finally pass to the limit in the integral
form of the differentiated equation on each compact strip.
All its spatial factors converge uniformly, and its coefficient
converges in $L^1$.  This proves local absolute continuity in time and
the asserted almost-everywhere equations.
\end{proof}

\begin{lemma}[Diffusion and square-martingale identity]
\label{lem:diffusion}
For $\gamma\in\mathcal U$, the equation
\begin{equation}
 \dd X_t=\gamma(t)u_\gamma(t,X_t)\,\dd t+\dd B_t,
 \qquad X_0=0,
 \label{eq:controlled-diffusion}
\end{equation}
has a unique strong solution on $[0,1]$.  On every strip $[0,T]$ with $T<1$,
\begin{equation}
 u_\gamma(t,X_t)=\int_0^t
                \Phi_{\gamma,xx}(s,X_s)\,\dd B_s.
 \label{eq:gradient-martingale}
\end{equation}
In particular, $q(t)=\E[u_\gamma(t,X_t)^2]$ belongs to $C^1[0,1)$,
$q(0)=0$, and
\begin{equation}
 q'(t)=\E[\Phi_{\gamma,xx}(t,X_t)^2].
 \label{eq:qprime}
\end{equation}
Under the approximations in Lemma~\ref{lem:regularity}, the coupled
diffusions converge uniformly on compact time strips, and $q_n$ and
$q_n'$ converge uniformly there.
\end{lemma}

\begin{proof}
The drift is measurable in time, bounded by $\gamma(t)$, and globally
Lipschitz in space on every compact time strip by
Lemma~\ref{lem:regularity}.  Picard iteration gives a unique strong
solution there.  These solutions agree on overlaps and extend to
time one because $\int_0^1\gamma<\infty$ and Brownian paths are
continuous.

For clarity, convergence needs no moment estimate for the
approximating paths.  Couple them using the same Brownian motion.
The drift difference is bounded by a common spatial Lipschitz
constant times $|X_n-X|$, plus
$|\gamma_n-\gamma|+\gamma_n\|u_n-u_\gamma\|_\infty$.
Gronwall's inequality gives a deterministic bound for
$\sup_{t\leq T}|X_n(t)-X(t)|$ tending to zero.  For step coefficients
and smooth terminal data, It\^o's formula on each step cancels the
drift of $u_n(t,X_n(t))$.  The step-boundary values agree.  Passing to
the limit in the stochastic integral by its $L^2$ isometry proves
\eqref{eq:gradient-martingale}.  Its initial value is zero by evenness.
The same isometry gives
$q(t)=\int_0^t\E[\Phi_{\gamma,xx}(s,X_s)^2]\,\dd s$.
The integrand is continuous by joint continuity, boundedness on
compact strips, and continuity of $X$.  This proves
\eqref{eq:qprime} and the stated convergence assertions.
\end{proof}

\begin{lemma}[Transfer on a constant interval]
\label{lem:constant-transfer}
Suppose $\gamma=c$ on an open interval $I\subset(0,1)$.
If $c>0$, every point of $I$ has a neighborhood on which one may
choose positive, bounded, nondecreasing finite step approximants
equal to $c$, with smooth terminal data and
$M_n\geq\|\gamma_n\|_\infty$.  At each time in this neighborhood,
all spatial derivatives and the corresponding diffusion
expectations of bounded polynomials in those derivatives converge.
Moreover, $\Phi_{\gamma,xx}>0$ throughout $I\times\R$.

For $c>0$, set $r=c\Phi_x$, $v=r_x$, $w=r_{xx}$, and $z=r_{xxx}$.
Then, within $I$,
\begin{align}
 (q')'(t)&=c^{-2}\E[w(t,X_t)^2-2v(t,X_t)^3],
 \label{eq:first-curvature}\\
 (q')''(t)&=c^{-2}\E[z(t,X_t)^2-12v(t,X_t)w(t,X_t)^2
                                      +6v(t,X_t)^4].
 \label{eq:second-curvature}
\end{align}
For $c=0$, $q'$ is strictly increasing on $I$.
\end{lemma}

\begin{proof}
To preserve a positive plateau, truncate and discretize the
coefficient separately to the left and right of a closed subinterval
of $I$, retain the value $c$ on that subinterval, and impose a
positive lower cutoff tending to zero and smaller than $c$.
All choices are monotone and converge in $L^1$.  Increase the terminal
smoothing parameter beyond the finite step bound.  The convergence
claims then follow from Lemmas~\ref{lem:regularity} and
\ref{lem:diffusion}.  The products involved are uniformly bounded,
and the spatial derivative bounds also give uniform continuity for
evaluating them on the coupled paths.

For strict positivity, choose $b\in I$ later than the time in
question and use \eqref{eq:constant-heat}.  The function
$h=\Phi_\gamma(b,\cdot)$ is smooth, convex, even, and differs from
$|x|$ by a bounded function.  Its derivative is therefore not constant.
Twice differentiating the logarithm in \eqref{eq:constant-heat}
gives
\[
 \Phi_{xx}(t,x)=\E_{t,x}^{\rm tilt}[h'']
                     +c\operatorname{Var}_{t,x}^{\rm tilt}(h').
\]
The tilted Gaussian law has a strictly positive density everywhere.
The variance is strictly positive and the first term is nonnegative.

Inside $I$ all functions are smooth in time as well as space, by
the heat formula.  Along $X$, the drifts of $r$, $v$, and $w$ are
respectively $0$, $-v^2$, and $-3vw$, and their diffusion coefficients
are $v$, $w$, and $z$.  These identities follow by differentiating
$r_t=-r_{xx}/2-rr_x$.  It\^o's formula for $v^2$ and then for
$w^2-2v^3$, with \eqref{eq:qprime}, proves
\eqref{eq:first-curvature}--\eqref{eq:second-curvature}.
All differentiations under the expectation are justified by the
bounded higher spatial derivatives on smaller intervals.

If $c=0$, the same argument gives
\[
 (q')'(t)=\E[\Phi_{xxx}(t,X_t)^2]>0.
\]
To see strictness, start the drift-free interval at any earlier
time in $I$.  Its Gaussian transition kernel implies that $X_t$ has
a strictly positive density on $\R$.  The function
$\Phi_{xxx}(t,\cdot)$ cannot vanish identically: otherwise
$\Phi(t,\cdot)$ would be an even quadratic polynomial, whereas it
is one-Lipschitz and differs from $|x|$ by a bounded function.
\end{proof}
